\documentclass[UTF8,11pt,a4paper]{ctexart}
\usepackage[a4paper,left=1.45cm,right=1.45cm,top=1.55cm,bottom=1.65cm]{geometry}
\usepackage{amsmath,amssymb,bm}
\usepackage{graphicx}
\usepackage{booktabs,tabularx,array,longtable}
\usepackage{enumitem}
\usepackage[dvipsnames,svgnames]{xcolor}
\usepackage{tikz}
\usetikzlibrary{arrows.meta,positioning,calc,fit,shapes.geometric}
\usepackage[most]{tcolorbox}
\usepackage{hyperref}
\usepackage{caption}
\usepackage{float}
\usepackage{fancyhdr}
\usepackage{titlesec}
\usepackage{microtype}
\IfFontExistsTF{Times New Roman}{\setmainfont{Times New Roman}}{\setmainfont{Liberation Serif}}
\IfFontExistsTF{Noto Sans CJK SC}{\setCJKmainfont{Noto Sans CJK SC}}{}
\IfFontExistsTF{Noto Sans CJK SC}{\setCJKsansfont{Noto Sans CJK SC}}{}
\IfFontExistsTF{DejaVu Sans Mono}{\setmonofont{DejaVu Sans Mono}}{\setmonofont{Latin Modern Mono}}
\definecolor{mainblue}{HTML}{0B5CAD}
\definecolor{softblue}{HTML}{EAF4FF}
\definecolor{darktext}{HTML}{1F2937}
\definecolor{softgray}{HTML}{F5F7FA}
\definecolor{accent}{HTML}{F59E0B}
\definecolor{goodgreen}{HTML}{2E7D32}
\definecolor{badred}{HTML}{B91C1C}
\definecolor{purple}{HTML}{6D5BD0}
\hypersetup{colorlinks=true,linkcolor=mainblue,urlcolor=DarkOrange,citecolor=mainblue}
\setlength{\parindent}{2em}
\setlength{\headheight}{14pt}
\setlength{\parskip}{0.22em}
\setlist[itemize]{leftmargin=2em,itemsep=0.23em,topsep=0.25em}
\setlist[enumerate]{leftmargin=2.2em,itemsep=0.23em,topsep=0.25em}
\renewcommand{\arraystretch}{1.2}
\captionsetup{font=small,labelfont=bf}
\titleformat{\section}{\Large\bfseries\sffamily\color{mainblue}}{\thesection}{0.7em}{}
\titleformat{\subsection}{\large\bfseries\sffamily\color{darktext}}{\thesubsection}{0.6em}{}
\titleformat{\subsubsection}{\normalsize\bfseries\sffamily\color{darktext}}{\thesubsubsection}{0.55em}{}
\pagestyle{fancy}
\fancyhf{}
\fancyhead[L]{\small\sffamily DSpark}
\fancyhead[R]{\small\sffamily \leftmark}
\fancyfoot[C]{\small\thepage}
\renewcommand{\headrulewidth}{0.4pt}
\newtcolorbox{infobox}[1]{enhanced,breakable,colback=softblue,colframe=mainblue!70!black,boxrule=0.8pt,arc=2mm,left=1.5mm,right=1.5mm,top=1mm,bottom=1mm,title=\bfseries #1}
\newtcolorbox{keybox}[1]{enhanced,breakable,colback=softgray,colframe=darktext!45,boxrule=0.6pt,arc=2mm,left=1.5mm,right=1.5mm,top=1mm,bottom=1mm,title=\bfseries #1}
\newtcolorbox{notebox}[1]{enhanced,breakable,colback=accent!8,colframe=accent!80!black,boxrule=0.7pt,arc=2mm,left=1.5mm,right=1.5mm,top=1mm,bottom=1mm,title=\bfseries #1}
\newcommand{\term}[1]{\textbf{\textcolor{mainblue}{#1}}}
\begin{document}

\section{论文基本信息}
\begin{center}
\begin{tabularx}{0.94\linewidth}{>{\bfseries}p{0.18\linewidth}X}
\toprule
题目 & DSpark\\
课件日期 & 2026-06-29\\
研究方向 & 大模型推理系统 / 并行调度\\
一句话概括 & 从系统层面对训练和推理链路进行并行化、调度与内存优化,以提高单用户生成速度和高并发有效吞吐。\\
\bottomrule
\end{tabularx}
\end{center}

\begin{infobox}{摘要要点}
从系统层面对训练和推理链路进行并行化、调度与内存优化,以提高单用户生成速度和高并发有效吞吐。 本说明按“研究问题—核心机制—基础公式—实验阅读—局限与优化”的顺序重新组织，不保留原始材料的封面、目录和页码对应。
\end{infobox}

\section{研究问题与动机}
这项工作的出发点不是简单地替换某个网络层，而是识别现有范式中最影响\term{质量、效率、可靠性或可部署性}的瓶颈。结合论文所讨论的问题，可以将研究动机概括为以下几点：
\begin{itemize}
\item 精度高但速度慢,抑或速度快但精度不足 这一矛盾构成投机解码的经典困境 Lau 博士云组会
\item 频繁出现搭配不当 如「of problem 」等不合理组合
\item 各位置token 独立预测 位置间缺乏依赖关系
\item 接受率呈断崖式下降 即所谓的「后缀衰减」现象
\item 并行草稿的关键问题:后缀衰减
\item 然而实际生产环境差异显著 不同任务特征各异 代码生成任务语法固定,8 个候选可全部接受 开放对话任务不确定性高,4 个候选可能全部失败
\end{itemize}

\begin{notebox}{理解重点}
阅读这类论文时，应把“问题是否真实存在”“方法是否直接解决该问题”“实验是否排除了其他解释”分开判断。仅凭更大的模型、更多数据或更长训练得到的提升，不能自动视为结构创新带来的收益。
\end{notebox}

\section{核心思想与整体流程}
从系统层面对训练和推理链路进行并行化、调度与内存优化,以提高单用户生成速度和高并发有效吞吐。 从信息流角度看，其基本流程可以整理为下图。

\begin{figure}[H]
\centering
\begin{tikzpicture}[>=Latex,node distance=0.78cm]
\tikzstyle{procstep}=[draw,rounded corners,fill=softblue,align=center,minimum width=3.05cm,minimum height=0.95cm,text width=2.75cm]
\node[procstep] (a) {请求与批次};
\node[procstep,right=of a] (b) {并行调度};
\node[procstep,right=of b] (c) {缓存/算子优化};
\node[procstep,right=of c] (d) {高速生成};
\draw[->,thick,mainblue] (a)--(b);
\draw[->,thick,mainblue] (b)--(c);
\draw[->,thick,mainblue] (c)--(d);
\end{tikzpicture}
\caption{核心信息流与处理阶段。}
\end{figure}

\subsection{方法组成与信息流}
\begin{itemize}
\item 公式4 :各位置最终概率分布 = 并行骨干基础输出+ 串行头偏移量 经归一化处理后获得 Lau 博士云组会
\item 设计思路:取长补短 1采用并行骨干确保生成速度 2叠加轻量串行头补偿预测精度
\item DSpark 的首要核心创新:半自回归生成架构
\item DSpark 采用全链路优化策略,同时提升三个关键维度 在速度、精度、调度三方面实现协同优化 Lau 博士云组会
\item DSpark 是一套全链路的系统级加速方案 整合了半自回归架构、置信度评分、硬件感知调度、在线校准四项关键技术
\item 公式1 :单token 实际延迟= ( 草稿生成耗时+ 验证耗时) / 接受token 数τ
\item DSpark 为行业提供了一个值得借鉴的工程实践范式 Lau 博士云组会
\end{itemize}

\begin{keybox}{结构解析}
\begin{itemize}
\item 是一套全链路的系统级加速方案 整合了半自回归架构、置信度评分、硬件感知调度、在线校准四项关键技术
\item 单用户速度提升60\\%-85\\% 高并发场景吞吐量提升约4 倍
\item 该方案已完成从论文研究到DeepSeek-V4 生产的完整落地
\item 单点技术红利收窄，全链路系统优化成为未来方向
\end{itemize}
\end{keybox}

下图补充展示模块连接、数据流或训练阶段之间的关系。
\begin{figure}[H]
\centering
\includegraphics[width=0.94\linewidth,height=0.43\textheight,keepaspectratio]{method_diagram.jpg}
\caption{模型结构、数据流或主要训练阶段。}
\end{figure}

\section{数学与机制理解}
本节给出理解该工作的基础数学结构。公式用于解释方法所属范式，并不替代原论文中的完整推导。
\begin{equation}
\text{Throughput}=\frac{\text{generated tokens}}{\text{time}},\qquad \text{latency}=t_{\mathrm{finish}}-t_{\mathrm{request}}
\end{equation}
系统工作必须同时考虑吞吐量、单请求时延、显存、并发公平性和实现复杂度。单一速度数字不足以完整评价真实部署收益。

\begin{keybox}{从公式回到方法}
应重点观察论文究竟改变了公式中的哪一部分：是输入表示、连接矩阵、条件变量、参数化方式、训练目标，还是求解与调度过程。真正的创新通常可以在这一层被准确定位。
\end{keybox}


\section{实验结果应如何阅读}
\begin{itemize}
\item 以Qwen3-14B 为例的关键数据 宏平均接受长度 较Eagle3 提升30.0\% 较DFlash 提升18.3\%
\item 在Qwen3 (4B/8B/14B) 与Gemma4-12B 上 全面超越Eagle3 与DFlash 基线
\item 覆盖3 大领域、9 个标准数据集 数学推理、代码生成、日常对话
\item 在所有测试场景中均表现稳定 Lau 博士云组会
\item 离线基准测试· 多模型全面验证
\item DeepSeek 将DSpark 部署于DeepSeek-V4 与现役MTP-1 基线进行真实流量对比
\item DSpark 平均接受长度 较Eagle3 提升26\%-31\% 较DFlash 提升16\%-18\%
\end{itemize}

\begin{notebox}{结果解读}
\begin{itemize}
\item Figure 7 帕累托前沿结果
\item DeepSeek 将 部署于DeepSeek-V4 与现役MTP-1 基线进行真实流量对比
\item 整条性能曲线位于基线外侧
\item 实际意义 在相同生成速度下，可承载更高总吞吐量 在相同并发量下，用户获得更快响应
\end{itemize}
\end{notebox}

下图用于直观呈现主要对比、消融、定性样例或效率结果。
\begin{figure}[H]
\centering
\includegraphics[width=0.94\linewidth,height=0.43\textheight,keepaspectratio]{result_evidence.jpg}
\caption{代表性实验对比、消融结果或定性样例。}
\end{figure}
实验部分不应只看单一最好数字。还需要核对模型规模、训练数据、训练步数、采样或解码预算、硬件条件以及是否使用额外蒸馏、检索或后处理。只有比较条件接近时，结果才能真正说明方法本身的优势。

\section{优点、局限与可优化空间}
\subsection{主要优点}
\begin{itemize}
\item \term{问题与方法对应明确}：论文围绕一个可识别的核心瓶颈展开，方法结构能够直接映射到该瓶颈。
\item \term{可复用的设计思想}：即使不完整复现模型，其中的分解、稀疏化、递归、条件控制或系统优化思想也可迁移到相邻任务。
\item \term{实验提供了多角度证据}：实验通常同时展示主结果、效率比较、案例或消融，使结论不完全依赖单一指标。
\end{itemize}

\subsection{局限与进一步优化}
\begin{itemize}
\item 需要进一步区分\term{结构收益}与更大算力、更多数据、不同训练技巧带来的收益，并在等参数、等计算预算下进行严格比较。
\item 对超参数、输入分布和模型规模的敏感性仍应系统分析；在一个数据集上的最优配置不一定能直接迁移。
\item 实际部署还要考虑显存峰值、并发、延迟抖动、失败案例和鲁棒性，而不仅是离线平均指标。
\item 可尝试将该方法与蒸馏、量化、缓存复用、动态路由或更强的表示学习结合，验证不同优化是否互补。
\end{itemize}

\section{结论}
总体而言，从系统层面对训练和推理链路进行并行化、调度与内存优化，以提高单用户生成速度和高并发有效吞吐。。 进一步需要注意：对大模型发展的启示 单点技术红利正在逐步收窄 真正的体验提升源于全链路的系统优化 每个环节的精细打磨将累积为显著的体验差异 Lau 博士云组会;DSpark 是一套全链路的系统级加速方案 整合了半自回归架构、置信度评分、硬件感知调度、在线校准四项关键技术。

\end{document}
